\section{Pre-training 的本质与局限}

\subsection{Pre-training 的独特优势}

Ilya 指出 pre-training 的两大核心优势：
\begin{enumerate}
    \item \textbf{数据量极其庞大}
    \item \textbf{不需要思考该训练什么数据}——答案是"everything"
\end{enumerate}

Pre-training 本质上是将"人类世界投射到文本上的全部内容"进行压缩学习。但它也非常难以推理——当模型犯错时，很难判断这是否因为某些内容在 pre-training 数据中支持不足。

\begin{warningbox}{Pre-training 没有人类类比}
Ilya 明确表示："I don't think there is a human analog to pre-training." 虽然有人将其类比为人类前 15 年的经历或生物进化，但他认为这些类比都不完美。Pre-training 数据量极其惊人，而人类只需极少数据就能达到更深刻的理解，且不会犯模型会犯的低级错误。
\end{warningbox}

\subsection{数据墙与下一步}

Pre-training 数据明确是有限的。当数据耗尽后，要么做某种"souped-up pre-training"（不同于以往的新配方），要么转向 RL，要么探索其他方向。无论如何，compute 已经很大了。

\begin{importantbox}{从 Scaling 时代回到 Research 时代}
\begin{itemize}
    \item 2012--2020：Research 时代——人们尝试各种想法，看什么有效
    \item 2020--2025：Scaling 时代——"scaling"这一个词吸走了房间里所有的空气，所有人做同样的事
    \item 2025+：\textbf{重回 Research 时代}——但配备了大型计算机
\end{itemize}
"Is the belief really that if you just 100x the scale, everything would be transformed? I don't think that's true."
\end{importantbox}

\subsection{本章小结}

Pre-training 是一个伟大但有限的配方。数据墙迫使领域重新思考训练方法。Ilya 认为我们正从 scaling 时代回到 research 时代——同样的创新驱动精神，但配备了数量级更大的计算资源。


